\begin{textsample}{POS Dim 3 – persona\_gemini – Score 18.00 – t686\_gemini.txt}  \label{ex:f3_pos_046}
The Paris Agreement ’s goal to limit global warming to 1.5°C demands radical changes across all \textbf{industries}.
While progress is being made in the \textbf{electricity} \textbf{sector}, where renewables now \textbf{supply} over 20 % of power in the EU, the \textbf{transport} \textbf{sector} is falling behind. \textbf{Emissions} from \textbf{transport} are actually rising, a \textbf{trend} fuelled by \textbf{increasing} \textbf{car} \textbf{sales}, especially of SUVs.
A \textbf{study} commissioned by Greenpeace and conducted by German \textbf{researchers} shows that the adoption of electric vehicles alone \textbf{will} not be sufficient to meet climate targets.
The research predicts that to stay within the \textbf{transport} \textbf{sector} 's carbon \textbf{budget}, a ban on the \textbf{sale} of new petrol and diesel \textbf{cars} must be implemented in Europe by 2028.
Norway serves as an example of how rapid change is possible, having \textbf{used} incentives to achieve fast \textbf{growth} in electric vehicle adoption with the aim for all new \textbf{cars} sold to be zero-emission by 2025.
However, a sustainable \textbf{transport} \textbf{system} requires broader \textbf{solutions}, such as investing in bike lanes, ensuring affordable \textbf{public} \textbf{transport}, and encouraging \textbf{car} sharing.

% matched lemmas: budget, car, electricity, emission, growth, increase, industry, public, researcher, sale, sector, solution, study, supply, system, transport, trend, use, will
\end{textsample}
